\documentclass[11pt,a4paper]{article}
\usepackage[margin=2.2cm]{geometry}
\usepackage[T1]{fontenc}
\usepackage[utf8]{inputenc}
\usepackage{lmodern}
\usepackage{xcolor}
\usepackage{hyperref}
\usepackage{listings}
\usepackage{booktabs}
\usepackage{longtable}
\usepackage{enumitem}
\usepackage{fancyhdr}
\usepackage{microtype}
\hypersetup{colorlinks=true,linkcolor=blue,urlcolor=blue}
\pagestyle{fancy}
\fancyhf{}
\lhead{Helm-DTN Standalone OCI Image Transfer}
\rhead{Implementation Report}
\cfoot{\thepage}
\definecolor{codebg}{RGB}{246,248,250}
\lstset{basicstyle=\ttfamily\small,backgroundcolor=\color{codebg},frame=single,
  breaklines=true,columns=fullflexible,numbers=left,numberstyle=\tiny,
  xleftmargin=1.5em,framexleftmargin=1em,showstringspaces=false}
\newcommand{\file}[1]{\texttt{#1}}
\newcommand{\lines}[1]{lines~#1}

\title{Extending Helm-DTN for Standalone OCI Image Transfer\\
\large Beginner-Oriented Design and Code Walkthrough}
\author{Project implementation documentation}
\date{5 August 2026}

\begin{document}
\maketitle
\begin{abstract}
This report explains how the original Helm-DTN project works, why its normal
chart-based pipeline cannot directly satisfy the requirement to provision a
remote registry before Kubernetes exists, and how the new Go implementation
adds standalone OCI-image transfer over the existing AAP2/\(\mu\)D3TN path.
Every important claim is connected to the current repository file and line
range. The intended endpoint is an OCI registry such as Harbor on the Space
side. Kubernetes is deliberately not required by the transfer stage; once the
registry contains the images, Kubernetes/containerd can pull them normally.
\end{abstract}

\tableofcontents
\newpage

\section{The requirement in plain language}
The system must move a real container image such as
\file{docker.io/library/nginx:alpine}, not merely a YAML reference to that
image. It must work without inventing a Helm chart, use the existing Go AAP2
client to cross the delay-tolerant network, place the received OCI content in
the configured Space-side Harbor/OCI registry, and verify that the destination
digest equals the source digest.

The boundary matters: this implementation ends at a usable registry. It does
not run \file{kubectl}, create a Pod, or assume that Kubernetes already exists.
That is necessary for provisioning use cases: Talos and Kubernetes images have
to arrive before the cluster can consume them.

\section{Helm and the original Helm-DTN code}
\subsection{What Helm actually does}
Helm is a Kubernetes package manager. A chart contains templates and values;
Helm renders these into Kubernetes objects and asks the Kubernetes API to
install them. An image field in a Deployment is only a string such as
\file{nginx:alpine}. Helm does not copy the image layers. After the Pod is
scheduled, the node's container runtime pulls those layers from a registry.

Therefore ``Helm cannot transfer a standalone image'' should be understood as
a responsibility boundary, not a defect in Helm. Helm manages releases and
manifests. OCI registries and container runtimes manage image manifests,
configurations, and layer blobs.

\subsection{How the supplied Helm-DTN pipeline is chart-coupled}
The original project adds image packaging around Helm charts:
\begin{enumerate}[leftmargin=*]
\item \file{pkg/helmdd/helmdd.go} renders a chart and discovers image strings
from generated Kubernetes manifests.
\item \file{pkg/helmdl/helmdl.go} downloads the discovered images.
\item \file{pkg/helmdp/helmdp.go} packages the chart and image archives for
transport.
\item \file{pkg/helmdtnt/receiver.go} receives the deployment, loads/pushes
images, and invokes Helm install or upgrade.
\end{enumerate}

The clearest evidence of the receiver's cluster dependency is
\file{pkg/helmdtnt/receiver.go}, \lines{170--182}, where the command type leads
to \file{installChart}, \file{upgradeChart}, or \file{uninstallChart}. The
actual Helm installation is at \lines{278--303}. At \lines{351--370}, Helm
initializes a Kubernetes action configuration and loads a chart from a local
path. This is correct for chart deployment, but it assumes both a chart and a
Kubernetes target.

The image loading part is similarly embedded in that deployment handler:
\file{pkg/helmdtnt/receiver.go}, \lines{212--239}, reads files from an
\file{Images} directory, constructs registry references, reads tarballs, and
pushes them. There is no independently addressable ``send this OCI image''
command.

\subsection{Why a fake chart is the wrong workaround}
A dummy chart could contain a Pod referring to an image, thereby causing the
existing detector to notice it. However, this would add an artificial chart,
trigger chart rewriting and Helm release behavior, require Kubernetes on the
receiver, and confuse registry provisioning with workload deployment. It also
makes failure recovery operate at the wrong abstraction: the desired unit is
an image, not a Helm release.

\section{New architecture}
The extension creates an explicit image path while retaining the existing AAP2
transport:

\begin{center}
Source registry $\rightarrow$ Go sender $\rightarrow$ OCI-layout TAR
$\rightarrow$ one AAP2 ADU $\rightarrow$ \(\mu\)D3TN
$\rightarrow$ Go receiver $\rightarrow$ Harbor/OCI registry.
\end{center}

There are three executables and two shared packages:
\begin{longtable}{p{0.31\textwidth}p{0.61\textwidth}}
\toprule
Component & Responsibility\\ \midrule
\file{cmd/image-sender/main.go} & Parse requests, export each selected platform,
frame it, and send it through the existing AAP2 client.\\
\file{cmd/image-receiver/main.go} & Subscribe on Space, stream the ADU to disk,
validate it, and push it to the configured registry.\\
\file{cmd/image-verify/main.go} & Independently compare source and destination
manifest digests.\\
\file{pkg/imageplan/plan.go} & Parse and validate command-line/image-list
source-to-target mappings.\\
\file{pkg/imagetransfer/transfer.go} & OCI export, payload framing, integrity
checks, safe extraction, registry push, and read-back verification.\\
\bottomrule
\end{longtable}

\section{Sender: exact code path and reasons}
\subsection{Inputs}
In \file{cmd/image-sender/main.go}, \lines{27--37}, repeated \file{-image} and
\file{-registry-target} options support one or many explicit mappings. The
\file{-images-file} option supports a reproducible list. The \file{-os} and
\file{-arch} options are needed because many registry tags point to a
multi-platform index; a provisioning transfer must select the platform that
the Space machines will execute.

At \lines{38--48}, the sender rejects missing network parameters and delegates
request construction to \file{imageplan.Build}. Centralizing parsing prevents
the sender and verifier from interpreting the same list differently.

\subsection{Temporary storage and AAP2 reuse}
At \lines{51--66}, the sender creates a temporary working directory when the
caller does not provide one and removes it on exit. Images may be hundreds of
megabytes, so using files prevents the entire payload from being retained as a
Go byte slice.

At \lines{67--74}, the code creates the project's existing
\file{AAP2TCPClient}, connects, and registers agent
\file{go-image-sender}. This is the key minimal-change decision: DTN routing,
framing, and compatibility continue to use the supplied AAP2 implementation.
No parallel DTN library or protocol was introduced.

\subsection{One image per ADU}
The loop at \lines{75--102} implements the transfer lifecycle. At \line{79},
\file{ExportOCI} pulls the image. At \lines{83--87}, the destination name is
stored in the header and the payload file is built. At \lines{88--94}, the
whole framed payload is submitted as one AAP2 Application Data Unit (ADU) and
the sender waits for the local acceptance response.

This code does \emph{not} implement application-level chunking. One image is
one ADU. Internally, Go's \file{io.Copy} streams bytes in bounded buffers, and
Bundle Protocol/\(\mu\)D3TN may fragment or transport data underneath, but
there is no numbered-chunk protocol or reassembly table in this application.
That distinction is important when diagnosing contact-capacity failures.

\section{OCI export and payload format}
\subsection{External libraries actually used}
The direct dependencies are recorded in \file{go.mod}. The relevant lines are:
\file{github.com/google/go-containerregistry v0.20.7} at \line{6} and
\file{github.com/google/uuid v1.6.0} at \line{7}.

\file{go-containerregistry} supplies registry authentication, reference
parsing, platform-aware pulling, OCI layout operations, and pushing. It avoids
a Docker daemon and avoids shelling out to Skopeo. \file{google/uuid} creates a
unique transfer identifier. Standard-library packages such as
\file{archive/tar}, \file{crypto/sha256}, \file{encoding/json}, and
\file{io} are not third-party libraries. Existing project dependencies provide
AAP2 protobuf messages and networking.

\subsection{Header and limits}
\file{pkg/imagetransfer/transfer.go}, \lines{25--38}, defines a 1 MiB maximum
header, a 64 GiB maximum archive, and the metadata fields: transfer ID, source,
target, creation time, source digest, archive SHA-256, and archive size.
Explicit bounds prevent malformed or hostile input from causing unbounded
allocation or disk consumption.

At \lines{40--55}, \file{ValidateTarget} rejects URLs, absolute paths, empty
components, and dot traversal. This is required because a sender-controlled
target is appended to the receiver's trusted registry prefix. Without this
check, input could attempt to escape that namespace or redirect a push.

At \lines{57--70}, \file{DefaultTarget} derives a destination repository and
tag when possible. Digest-only source references require an explicit target
because a destination push needs a usable repository name and tag/digest
mapping rather than an ambiguous default.

\subsection{Pulling a real image into an OCI layout}
At \lines{72--109}, \file{ExportOCI} parses the source and calls
\file{remote.Image} with keychain authentication and the selected platform
(\lines{75--85}). It creates an OCI image layout at \lines{88--99}, appends the
image, turns the layout directory into a TAR at \line{100}, and hashes the TAR
at \lines{104--108}. The result is a standards-based OCI object containing the
manifest, configuration, and layer blobs---not a text reference.

\subsection{Wire framing}
At \lines{111--150}, \file{CreatePayload} writes exactly:
\begin{verbatim}
[8-byte unsigned big-endian JSON length][JSON Header][OCI-layout TAR bytes]
\end{verbatim}
The length prefix lets the receiver know where metadata ends and binary image
data begins. The JSON is readable and versionable; the TAR remains streamable.
At \line{142}, \file{io.Copy} copies file-to-file rather than loading the
archive into memory.

\subsection{Integrity checks on Space}
At \file{pkg/imagetransfer/transfer.go}, \lines{152--205},
\file{DecodePayload} reads the bounded header exactly, checks all required
fields, validates the target and archive size, extracts only the declared
number of bytes, and recomputes size and SHA-256. Corruption or truncation is
rejected before any registry push.

The source manifest digest and archive checksum serve different purposes. The
archive checksum proves transport-byte integrity. The manifest digest proves
that the registry-visible image content is the same image, even though registry
storage need not preserve the surrounding TAR byte-for-byte.

\section{AAP2 client modification}
The extension intentionally makes a small change to
\file{pkg/aap2_client/aap2_sync_client.go}. The existing \file{SendFile} path
at \lines{238--260} is reused. Its new size check is at \lines{245--251};
streaming to the connection remains at \lines{252--258}. This avoids announcing
one length and streaming another.

The in-memory receive paths use \file{io.ReadFull} at \lines{78--85} and
\lines{162--168}, ensuring a short TCP read is not mistakenly treated as a
complete payload. The new \file{ReceiveAduFile} method at \lines{171--197}
streams exactly \file{PayloadLength} bytes to a file using \file{io.CopyN} at
\line{188}. This is necessary because normal network reads are allowed to
return fewer bytes than requested and because provisioning images should not
require an equally large RAM allocation.

These are transport-hardening changes, not a replacement for AAP2. The sender
and receiver still configure, send messages, send response statuses, and carry
ADUs through the original Go client.

\section{Receiver: from DTN to Harbor}
In \file{cmd/image-receiver/main.go}, \lines{17--24} define the Space AAP2
address, subscriber agent ID, optional secret, registry prefix, work directory,
and a testing-only insecure-registry switch. Harbor credentials are read at
\lines{29--40} from \file{REGISTRY_CREDS}; \file{HARBOR_CREDS} is supported for
backward compatibility. No value means anonymous registry access.

The outer loop at \lines{44--49} reconnects after an AAP2 failure. DTN links
and long-running nodes can be intermittent, so a receiver must not exit forever
on one lost TCP connection.

At \lines{52--65}, the receiver uses the existing AAP2 TCP client and
subscribes as \file{image-transfer}. At \lines{66--87}, it receives the ADU
metadata, acknowledges the AAP2 message, streams the declared payload into a
unique file, processes it, and removes the temporary payload.

At \lines{90--104}, \file{handle} decodes and verifies the payload, combines
the administrator-controlled registry prefix with the validated relative
target, pushes the OCI image, and logs its verified digest. There is no Helm
install and no Kubernetes client in this executable.

\section{Registry push and end-to-end verification}
\file{PushOCI} in \file{pkg/imagetransfer/transfer.go}, \lines{208--260}, safely expands the OCI
layout, opens the image, parses the destination reference, selects basic or
anonymous authentication, and writes it through the Registry V2 API. After the
write it calls \file{remote.Head} and compares the registry's reported digest
with the local image digest. A successful HTTP upload alone is insufficient;
the read-back proves that the expected manifest is addressable at the target.

The separate program \file{cmd/image-verify/main.go} provides independent
acceptance evidence. At \lines{28--33} it loads the same validated mapping and
credentials. At \lines{34--54} it calculates the selected-platform source
digest, calculates the remote target digest, compares them, and prints
\file{OK} only on equality. Any failure causes a non-zero process exit at
\lines{55--57}. The digest helper at \lines{75--100} deliberately uses source
keychain authentication and destination registry credentials separately.

\section{Image-list methodology}
\file{pkg/imageplan/plan.go}, \lines{17--98}, accepts either command-line pairs or a text file.
Each non-comment line is either \file{SOURCE} or \file{SOURCE=TARGET}. It
rejects missing sources, mismatched repeated options, malformed mappings,
unsafe targets, and an empty final plan. This shared package makes batch
transfers reproducible and keeps validation identical between transfer and
verification.

Two examples are supplied. \file{examples/kubernetes-images.txt} is a small
demonstration. \file{examples/talos-v1.13.6-kubernetes-v1.36.2-images.txt} is
a pinned provisioning set. Pinning versions is necessary because a mutable
``latest'' set makes results and future incident analysis irreproducible.

\section{Local acceptance script and DTN capacity}
\file{scripts/run-local-acceptance.sh}, \lines{4--11}, selects the image list,
registry, Earth and Space AAP2 endpoints, destination EID, contact duration,
and contact rate. At \lines{19--23}, it starts the existing test containers and
configures the Earth-to-Space contact. The default rate is 10,000,000 bytes/s
for 3,600 seconds, giving about 36 GB of scheduled capacity.

This setting was needed because the earlier 100,000 bytes/s contact offered
only about 360 MB total. Several Talos images individually approach or exceed
that practical budget once bundle overhead and multiple images are considered;
\(\mu\)D3TN consequently reported insufficient contact payload capacity. This
was a DTN scheduling capacity issue, not an OCI or Harbor failure.

At \lines{25--34}, the script starts the receiver and sender. At
\lines{36--48}, it waits for the expected number of registry pushes rather than
assuming local AAP2 acceptance means remote delivery. Finally, at
\lines{50--55}, it runs independent digest verification and reports completion.

\section{Security, reliability, and limitations}
\begin{itemize}[leftmargin=*]
\item Integrity is checked with archive SHA-256 before pushing and manifest
digest after pushing. This detects corruption; it is not a digital signature.
For hostile environments, add signed transfer metadata or verify Cosign/OCI
signatures against a trusted policy.
\item \file{-insecure-registry} permits HTTP and must remain limited to local
testing. Real Harbor should use TLS and appropriately scoped robot credentials.
\item Credentials are currently passed as \file{user:password} in an
environment variable. A colon in the username is unsupported, and environment
visibility may be undesirable. A secret file or credential helper is a future
hardening improvement.
\item One image equals one ADU. There is no application retry ledger,
resumable-chunk protocol, deduplication across images, or garbage-collection
policy for failed receiver files.
\item Sender acceptance confirms that local \(\mu\)D3TN accepted the ADU; it
does not itself confirm remote Harbor completion. Receiver logs and
\file{image-verify} provide that evidence.
\item A local Registry V2 endpoint has been used for functional validation.
The same API is used by Harbor, but final deployment should still test actual
Harbor TLS, project permissions, quotas, and retention rules.
\end{itemize}

\section{What was modified}
\begin{longtable}{p{0.36\textwidth}p{0.56\textwidth}}
\toprule
File & Modification and reason\\ \midrule
\file{cmd/image-sender/main.go} & New standalone sender; removes the chart
prerequisite and emits one image per AAP2 ADU.\\
\file{cmd/image-receiver/main.go} & New registry-only Space receiver; removes
the Kubernetes/Helm requirement from image provisioning.\\
\file{cmd/image-verify/main.go} & New independent digest verifier; distinguishes
local acceptance from completed remote delivery.\\
\file{pkg/imagetransfer/transfer.go} & New OCI and payload implementation with
bounds, checksums, safe extraction, registry push, and read-back.\\
\file{pkg/imageplan/plan.go} & New shared, validated batch mapping parser.\\
\file{pkg/aap2_client/aap2_sync_client.go} & Hardened exact-length I/O and added
disk-streamed receiving for large image ADUs.\\
\file{pkg/*/*_test.go} & Added tests for malformed mappings, framing,
integrity, path safety, and exact AAP2 reads.\\
\file{examples/*.txt} & Added reproducible Kubernetes and Talos image plans.\\
\file{scripts/run-local-acceptance.sh} & Added end-to-end local DTN, registry,
count, and digest acceptance workflow.\\
\file{STANDALONE-IMAGES.md}, \file{PROVENANCE.md}, \file{README.md} & Added usage,
scope, dependency provenance, and project navigation documentation.\\
\bottomrule
\end{longtable}

\section{How to build and run}
Build the three programs from the repository root:
\begin{lstlisting}[language=bash,numbers=none]
go build -o bin/image-transfer-send ./cmd/image-sender
go build -o bin/image-transfer-receive ./cmd/image-receiver
go build -o bin/image-transfer-verify ./cmd/image-verify
\end{lstlisting}

On Space, start the receiver (use TLS in production):
\begin{lstlisting}[language=bash,numbers=none]
export REGISTRY_CREDS='robot-user:password'
bin/image-transfer-receive \
  -tcp 127.0.0.1:4242 \
  -registry harbor.space.example/provisioning
\end{lstlisting}

On Earth, send one real image:
\begin{lstlisting}[language=bash,numbers=none]
bin/image-transfer-send \
  -tcp 127.0.0.1:4242 \
  -destination-eid dtn://space.dtn/image-transfer \
  -image docker.io/library/nginx:alpine \
  -registry-target library/nginx:alpine \
  -os linux -arch amd64
\end{lstlisting}

After delivery, verify the batch with \file{image-transfer-verify}. The
supervisor's requested implementation ends when the standalone OCI image is
present in the Space-side Harbor/OCI registry and its digest has been verified.
This project does not ask Kubernetes to pull the image and does not start a
container or Pod.

\section{Testing status and honest completion boundary}
Unit tests cover the newly introduced parsing, framing, integrity, target
validation, and exact AAP2 file I/O. A real
\file{registry.k8s.io/pause:3.10} image was transferred through the local DTN
test environment into a Registry V2 destination and verified at digest
\file{sha256:7c38f24774e3cbd906d2d33c38354ccf787635581c122965132c9bd309754d4a}.

The full pinned 17-image Talos acceptance run has not yet completed in one
uninterrupted execution, and a real Space Harbor has not yet been exercised.
Consequently, the implementation is present and the core path is demonstrated,
but final acceptance should require: \file{go test ./...}, one uninterrupted
17/17 run, independent 17/17 digest verification, and a deployment-environment
Harbor test.

\section{Beginner's mental model}
Keep four acknowledgements separate:
\begin{enumerate}[leftmargin=*]
\item The source registry returned a valid image and digest.
\item Earth \(\mu\)D3TN accepted the image ADU locally.
\item Space received, checksum-validated, and pushed the OCI image.
\item A fresh registry read returned the same manifest digest.
\end{enumerate}
The fourth statement establishes the supervisor's intended final result. Image
consumption by Kubernetes is explicitly outside this implementation.

\end{document}
